\section{Introduction}
\label{intro:section}

The Weil classes are a particularly concrete test of the Hodge conjecture:
they form a rational two-dimensional space defined by an imaginary
quadratic action, while their algebraic representatives generally do not
come from products of divisor classes.  This paper establishes their
algebraicity in the split eight-dimensional case.

\subsection{Statement and scope}

Let $K=\Q(\sqrt{-d})$, where $d$ is a positive squarefree integer, and let
$A$ be a complex abelian variety of dimension eight.  Suppose that
\[
 \iota:K\hookrightarrow\End(A)\otimes_\Z\Q
\]
is a $\Q$-algebra embedding and that a polarization $\lambda$ has Rosati
involution restricting to complex conjugation on $\iota(K)$.  Write
$\sigma,\bar\sigma:K\hookrightarrow\C$ for its two embeddings.
The \emph{Weil-type condition} is
\[
 \dim_\C H^{1,0}(A)_\sigma=
 \dim_\C H^{1,0}(A)_{\bar\sigma}=4.
\]
On the $K$-vector space $V=H_1(A,\Q)$, of dimension eight, let $\psi$
be the rational Riemann form of $\lambda$ and let $J$ denote the action
of $\iota(\sqrt{-d})$.  The associated Hermitian form, linear in its first
variable, is
\begin{equation}
 h(x,y)=\psi(Jx,y)+\sqrt{-d}\,\psi(x,y).
 \label{intro:hermitian}
\end{equation}
We say that the tuple is \emph{split} if $V$ contains a four-dimensional
$K$-subspace on which $h$ vanishes identically.  This definition fixes the
convention without referring to a choice of sign for the discriminant.

The rational Weil space is
\[
 W_K(A)=\bigwedge\nolimits_K^8 H^1(A,\Q)\subset H^8(A,\Q).
\]
Its embedding is specified by the decomposition
\[
 W_K(A)\otimes_\Q\C=
 \bigwedge\nolimits_\C^8H^1(A,\C)_\sigma
 \oplus
 \bigwedge\nolimits_\C^8H^1(A,\C)_{\bar\sigma}.
\]
Thus $\dim_\Q W_K(A)=2$, and the Weil-type condition places it in
$H^{4,4}(A)$.

\begin{theorem}
\label{intro:main}
For every split Weil-type tuple $(K,A,\iota,\lambda)$ as above,
\[
 W_K(A)\subset
 \im\bigl(\CH^4(A)\otimes_\Z\Q
                \xrightarrow{\ \mathrm{cl}\ }H^8(A,\Q)\bigr).
\]
The assertion holds for every imaginary quadratic field, every compatible
polarization type, and every member of the split Weil-type families,
including members with additional endomorphisms.
\end{theorem}

The conclusion concerns rational linear combinations of algebraic cycles.
It resolves the Weil-class algebraicity conjecture for these split
eightfolds.  It makes no assertion about all Hodge classes on arbitrary
abelian eightfolds or about the Hodge conjecture for arbitrary smooth
projective varieties.

\subsection{Historical setting}

Weil introduced the exceptional determinant classes in 1977
\cite{weil}.  Markman's work proves their algebraicity on all abelian
fourfolds of Weil type and on split abelian sixfolds of Weil type;
it also yields the rational Hodge conjecture for abelian varieties of
dimension at most five \cite{markman-sixfold,markman-survey}.
The split eightfold case is a further instance beyond those results.
The February 2026 version of
\cite[Question~11.4 and Section~12]{markman-survey} proposes an eightfold
construction using a weaker semiregularity condition.  Both verification
of that condition for the example and its sufficiency for the required
deformation theorem remain open there.

Our proof proceeds through a symplectic construction and modules over
theta section rings.  Its main additional ingredients are a Lagrangian
with prescribed homology and cohomological detection properties, a finite
construction of weighted Floer operations and their obstruction identity,
and an alternating specialization argument for periodic module
cohomology.  The last argument transports algebraic Chern characters
without requiring a deformation theorem for a mirror object.

\subsection{Proof architecture}

The proof has four stages.  First, Section~\ref{hodge:section} identifies
all rational split Hermitian spaces with a fixed model and constructs
a middle-dimensional class $\alpha$.  Scalar graph tori pair trivially with
the exceptional summand of $\alpha$, while one further graph detects it.
Section~\ref{top:section} realizes a positive multiple
of $\alpha$ by an embedded spin Lagrangian $L$, graded modulo four.  The
maps $H_2(L)\to H_2(T)$ and $H_6(L)\to H_6(T)$ are injective, and a
specified space of ambient two-forms restricts trivially in cohomology.

Second, Sections~\ref{local:section}--\ref{cat:section} construct the
weighted Floer module.  The analytic construction works with finite
energy and arity cutoffs.  It includes boundary gluing, compatible local
thickenings, integration on overlaps, and the disk identity with one
interior marked point.  The two homology injections force the
cohomological obstructions to vanish.  A residue field of an ultraproduct
then turns the finite constructions into exact algebraic identities over
a power-series ring $R=\kappa[[t]]$.

Third, Section~\ref{theta:section} computes the directed graph algebra
as a theta section ring $\mathscr S$.  Its relative projective spectrum
$\mathscr X=\Proj\mathscr S$ is smooth and projective over $R$.
At $t=0$, Sections~\ref{mirror:section} and~\ref{per:section} apply
ordinary mirror symmetry to finite lists of graph branes and recover
the new module as a perfect object in a category with period-four
grading.  Its coherent cohomology sheaves give a finitely generated
graded module in each residue degree.  Euler pairings with the scalar
graphs and the extra graph show that its alternating Chern character
lies outside the span of the powers of the polarization.

Finally, Section~\ref{lift:section} lifts finite generation and the
alternating character across $R$.  The $t$-torsion terms cancel in the
alternating sum of derived restrictions.  The theta ring can be realized
at complex periods avoiding specified countably many analytic loci.
Section~\ref{spread:section} chooses such a period, isolates a nonzero
Weil component, and spreads its marked class to the entire family using
proper Hilbert parameter spaces.  The quadratic-field action supplies
a second independent Weil class; the initial isogeny reduction then
proves Theorem~\ref{intro:main}.

\subsection{Conventions and inputs}

The complex abelian eightfolds $A_\Pi$ and the real symplectic
sixteen-torus $T$ play different roles.  The symplectic form on $T$ is
fixed throughout.  Variation of $\Pi$ enters the Floer operations through
an additional formal relative two-class.  Only the theta multiplication
constants are prescribed under the field embeddings used to realize a
complex period; there is no assertion that the new Lagrangian's Floer
series converge at that period.

We use established immersion and embedding results, the ordinary
family Floer theorems, and standard algebraic and analytic geometry with
their hypotheses stated at the point of application.  The finite
weighted virtual construction and its comparison identities are proved
in Sections~\ref{local:section}--\ref{cat:section}.  In particular, the
ordinary mirror theorems are applied only to tautologically unobstructed
graphs.  The new brane enters through its module and the proof does not
assume a mirror equivalence for the deformed bounding-cochain category.

Homology and cohomology are rational unless another coefficient field is
displayed.  Periodic degrees are indexed by $\Z/4$; the function
$r\mapsto(-1)^r$ is well defined on this group.  Chern characters are
rational algebraic Chern characters, with their Betti or $\ell$-adic
cycle classes specified where needed.
